\documentclass[12pt]{article}
\usepackage[T2A]{fontenc}
\usepackage[utf8]{inputenc}
\usepackage[russian]{babel}
\usepackage{amsmath,amssymb}
\usepackage{geometry}
\geometry{margin=2.5cm}
\usepackage{indentfirst}

\title{Квантовые поправки к гравитационному линзированию\\ в присутствии аксионов}
\author{С.~Т.~Белов, Х.~ван дер Берг}
\date{2024}

\begin{document}
\maketitle

\section*{Аннотация}
Рассмотрено распространение света в пространстве-времени с аксионным полем, связанным с электромагнитным дуальным тензором. Вычислены однопетлевые квантовые поправки к углу отклонения луча и времени задержки. Показано, что эффект может быть заметен для линз с прицельным параметром порядка радиуса горизонта.

\vspace{0.5cm}
\noindent\textbf{Ключевые слова:} аксионы; гравитационное линзирование; квантовые поправки; дуальный тензор; космология.

\vspace{0.3cm}
\noindent\textbf{Раздел журнала:} Физика. \quad \textbf{Файл:} 1 из 3.

\section{Введение}
Аксионы — гипотетические частицы, взаимодействующие с электромагнитным полем через дуальный тензор. Мы исследуем их влияние на гравитационное линзирование.

\section{Квантовые поправки}
Однопетлевая поправка к углу отклонения имеет вид
\begin{equation}
\delta \alpha = \frac{g_{a\gamma\gamma}^2}{16\pi^2} \frac{m_a^2}{b^2} \ln\left(\frac{\Lambda^2}{m_a^2}\right),
\end{equation}
где $g_{a\gamma\gamma}$ — константа связи, $m_a$ — масса аксиона, $b$ — прицельный параметр.

\section{Заключение}
Обсуждаются наблюдательные следствия для будущих миссий.

\end{document}